\begin{proofbox}
	\textbf{\textcolor{CProof}{思路提示}}
	数列连续三项成等差，意味着相邻两项的差相等，从而得到中间项与前后项的关系。

	\medskip
	\textbf{\textcolor{CProof}{正式推导}}
	\begin{description}[style=nextline, leftmargin=2.8em, labelsep=0.5em, itemsep=0.55em]
		\item[\textbf{步骤一（设公差）：}]
			设 $a_{n-1},a_n,a_{n+1}$ 成等差数列，公差为 $d$，则
			\[
				a_n - a_{n-1} = d,\quad a_{n+1} - a_n = d.
			\]
			\par\textbf{理由：} 等差数列定义：任何相邻两项的差相等。
		\item[\textbf{步骤二（等量传递）：}]
			由两个差都等于 $d$，得
			\[
				a_n - a_{n-1} = a_{n+1} - a_n.
			\]
			\par\textbf{理由：} 等号传递性：$a_n-a_{n-1}=d$，$a_{n+1}-a_n=d$，故二者相等。
		\item[\textbf{步骤三（移项合并）：}]
			将上式中的两项移到同一侧：
			\[
				a_n - a_{n-1} - (a_{n+1} - a_n) = 0,
			\]
			整理得
			\[
				a_n - a_{n-1} - a_{n+1} + a_n = 0,
			\]
			即
			\[
				2a_n - a_{n-1} - a_{n+1} = 0.
			\]
			\par\textbf{理由：} 去括号并合并同类项。
		\item[\textbf{步骤四（最终变形）：}]
			移项得
			\[
				2a_n = a_{n-1} + a_{n+1},
			\]
			两边除以 $2$ 即得算术平均形式：
			\[
				a_n = \frac{a_{n-1}+a_{n+1}}{2}.
			\]
			\par\textbf{理由：} 代数恒等变形。
	\end{description}

	\medskip
	\textbf{\textcolor{CProof}{结论回扣}}
	\textbf{由以上推导可见，只要连续三项成等差数列，中间项一定等于前后两项的算术平均数。该结论既是等差数列的性质，也为求项和证明等差提供了工具。}
\end{proofbox}
